\documentclass[12pt,a4paper]{article}


% Package declarations for formatting, math, and code listings
\usepackage[utf8]{inputenc}
\usepackage{geometry}
\geometry{a4paper, margin=1in}
\usepackage{amsmath}
\usepackage{graphicx}
\usepackage{booktabs}
\usepackage{listings}
\usepackage{xcolor}
\usepackage{hyperref}

% Code snippet styling configurations
\definecolor{mGreen}{rgb}{1,0,0}
\definecolor{mGray}{rgb}{0.5,0.5,0.5}
\definecolor{mPurple}{rgb}{0.58,0,0.82}
\definecolor{backgroundColour}{rgb}{0.95,0.95,0.92}

\lstset{
    backgroundcolor=\color{backgroundColour},   
    commentstyle=\color{mGreen},
    keywordstyle=\color{green},
    numberstyle=\tiny\color{mGray},
    stringstyle=\color{mPurple},
    basicstyle=\ttfamily\footnotesize,
    breakatwhitespace=false,         
    breaklines=true,                 
    captionpos=b,                    
    keepspaces=true,                 
    numbers=left,                    
    numbersep=5pt,                  
    showspaces=false,                
    showstringspaces=false,
    showtabs=false,                  
    tabsize=2,
    language=C
}

% Document Metadata
\title{\textbf{Technical Report on Cinema Ticket Booking System}}
\author{GROUP F MEMBERS\and
RUGABA BENEDICT-25/U/BIE/01417/PE\and
MUKINE JOEL-22/U/BIO/1688/PD\and
AUGUSTINE CEASAR OLOYA-24/U/BIE/03944/PE\and
NAIWUMBWE ASSUMPTA-24/U/BIE/08557/PE
}
\date{\today}

\begin{document}

\maketitle
\newpage
\tableofcontents
\newpage

\section{Introduction}
This report details the architectural design and technical implementation of a lightweight, terminal-based Cinema Ticket Booking System written in the C programming language. The application acts as a Point of Sale environment managing inventory, client reservations, dynamic tiered pricing structures, and financial auditing metrics.

\section{System Build up and Data Structures}
The system maintains data state utilizing two main \textit{struct} definitions to control memory allocation arrays without running into pointer or memory leak issues.

\subsection{Movie Section}
The \textit{Movie} structure manages the current physical seating state, inventory caps, and individual revenue tracking.
\begin{lstlisting}[caption=Movie Struct Definition]
typedef struct {
    int id;
    char title[50];
    int seats[20]; // 0 = Available, 1 = Booked
    int tickets_sold;
    float total_sales;
} Movie;
\end{lstlisting}

\subsection{Booking Section}
The \textit{Booking} structure archives completed client transaction parameters. to be specific, the program obtains the name, the seat they want and other different parameters.
\begin{lstlisting}[caption=Booking Struct Definition]
typedef struct {
    int booking_id;
    char customer_name[50];
    int movie_id;
    int seat_number;
    char customer_type; // 'C' = Child, 'S' = Student, 'A' = Adult
    int is_weekend;     // 1 = Yes, 0 = No
    float final_price;
} Booking;
\end{lstlisting}

\section{The Logic}
Here we shall see the logic behind the source code.

\subsection{Pricing Parameters}
Let $P_{base}$ represent the base cost depending on the customer age type $T$:
\begin{equation}
P_{base} = 
\begin{cases}
\$5.00 & \texttt{if } T = \text{'C' (Child)} \\
\$8.00 & \texttt{if } T = \text{'S' (Student)} \\
\$12.00 & \texttt{if } T = \text{'A' (Adult)} 
\end{cases}
\end{equation}

\subsection{Surcharge and Group Discounts}
Let $W$ be the weekend modifier ($W = 0.20$ if weekend, $0$ otherwise), and $D$ be the group bulk discount modifier applied if total grouped tickets $N \ge 3$ ($D = 0.10$, $0$ otherwise). The formula for final single-ticket evaluation $P_{final}$ is modeled as follows:
\begin{equation}
P_{final} = (P_{base} + (P_{base} \times W)) \times (1 - D)
\end{equation}

\section{Functional Matrix Breakdown}
The application exposes a structured pipeline for running administrative workflows.

\begin{table}[h!]
\centering
\caption{System Operations and Runtime Status}
\begin{tabular}{lll}
\toprule
\textbf{Function Call} & \textbf{Operational Scope} & \textbf{Implementation Status} \\
\midrule
\texttt{init\_system()} & Bootstraps core catalogs & Complete \\
\texttt{book\_tickets()} & Registers client data & Complete \\
\texttt{search\_booking()} & Iterates via client matching tags & Complete \\
\texttt{update\_booking()} & Swaps seat allocations safely & Complete \\
\texttt{cancel\_booking()} & Purges records using ID masks & Complete \\
\texttt{display\_available\_seats()} & Renders matrix layout & Incomplete / Pending \\
\bottomrule
\end{tabular}
\label{tab:functions}
\end{table}

\section{Conclusion and Future Work}
In conculsion, the program above depicits the basic cinema ticket booking
\begin{itemize}
    \item File I/O serialization to persist memory tables across session resets.
    \item Fixed size buffer overflow shields (\texttt{\%49[\^{}\textbackslash n]} rules) for input capture safety.
\end{itemize}

\end{document}
